\chapter{Decomposição em Sub-QUBOs}
\label{chap:decomposicao-subqubo}

Após a decomposição adaptativa, o problema deixa de ser expresso em termos de
processos individuais e passa a ser definido sobre um conjunto de entidades. Cada entidade pode corresponder a um processo isolado ou a um
\textit{bundle}. Seja

\begin{equation}
    \mathcal{E} = \{e_0, e_1, \dots, e_{N-1}\}
\end{equation}
\listeq{Conjunto de entidades após a decomposição adaptativa}

o conjunto de entidades resultante, onde cada entidade $e_i$ tem peso de
\acs{CPU} $w_i$ e memória residente $r_i$. 

Para $K$ núcleos, a formulação
continua a exigir $n_q = N \cdot K$ qubits lógicos. Assim, mesmo depois do agrupamento em \textit{bundles}, o problema
pode continuar a exceder o limite configurado $q_{max}$. Quando
$N \cdot K \leq q_{max}$, a plataforma resolve um único \acs{QUBO}. Quando
$N \cdot K > q_{max}$, é usada uma pipeline iterativa baseada em
\textit{sub-QUBOs}.

\section{Construção do QUBO Global}
\label{sec:construcao-qubo-global-subqubo}

O primeiro passo da pipeline iterativa consiste em construir uma matriz
\acs{QUBO} global $Q^{G}$ sobre todas as entidades de $\mathcal{E}$. Esta matriz
é construída com a mesma formulação apresentada anteriormente.

Para cada entidade $e_i$ e núcleo $k$, a variável binária correspondente é:

\begin{equation}
    x_{ik} =
    \begin{cases}
        1, & \text{se a entidade } e_i \text{ for atribuída ao núcleo } k,\\
        0, & \text{caso contrário.}
    \end{cases}
\end{equation}
\listeq{Variável de atribuição entidade-núcleo na decomposição iterativa}

No entanto, quando o número de variáveis excede
$q_{max}$, esta matriz não é enviada diretamente para o \textit{solver}. Em vez
disso, é usada como referência para construir sucessivos \textit{sub-QUBOs}.

\section{Extração de um Sub-QUBO}
\label{sec:extracao-subqubo}

As entidades construidas na decomposição adaptiva são divididas em grupos

\begin{equation}
    \mathcal{G}_0, \mathcal{G}_1, \dots, \mathcal{G}_{T-1},
\end{equation}
\listeq{Partição das entidades em grupos}

onde cada grupo contém no máximo

\begin{equation}
    m = \left\lfloor \frac{q_{max}}{K} \right\rfloor
\end{equation}
\listeq{Número máximo de entidades por sub-QUBO}

entidades. 

Estes grupos são feitos através de duas estratégias:

\begin{itemize}
    \item \texttt{WEIGHT\_DESCENDING}: ordena os \textit{bundles} por carga
    de CPU decrescente;

    \item \texttt{COUPLING\_DESCENDING}: começa pelo \textit{bundle} com
    maior magnitude de carga e adiciona os \textit{bundles} com maior
    acoplamento, dado por $2W_iW_j$, aos que já pertencem ao grupo. Desta forma, as interações mais fortes permancem dentro do mesmo \textit{sub-QUBO}.
\end{itemize}

Para cada grupo $\mathcal{G}_t$, é extraída de $Q^{G}$ a submatriz associada às
variáveis das entidades desse grupo. 

Este processo preserva os termos internos entre entidades do mesmo grupo. No
entanto, os termos com entidades já resolvidas em iterações anteriores
deixam de aparecer, porque essas entidades já têm uma
atribuição fixa. Para incorporar o seu efeito, a implementação mantém um vetor
de carga acumulada por núcleo:

\begin{equation}
    \boldsymbol{\phi}^{(t)} =
    \left(\phi^{(t)}_0, \phi^{(t)}_1, \dots, \phi^{(t)}_{K-1}\right),
\end{equation}
\listeq{Carga acumulada por núcleo antes da iteração \(t\)}

inicializado com a carga dos processos de tempo real $\mathcal{P}_{RT}$. Para uma entidade
$e_i$ do grupo atual e um núcleo $k$, a diagonal do \textit{sub-QUBO} associado à atribuição
$(e_i,k)$ é corrigida
por:

\begin{equation}
    Q^{(t)}_{(e_i,k),(e_i,k)}
    \leftarrow
    Q^{(t)}_{(e_i,k),(e_i,k)}
    +
    2w_i\phi^{(t)}_k.
\end{equation}
\listeq{Propagação da carga acumulada como bias diagonal}

Este termo vem da expansão do balanceamento de carga quando já existe carga fixa
no núcleo $k$:

\begin{equation}
    \left(\phi^{(t)}_k + \sum_{e_i \in \mathcal{G}_t} w_i x_{ik} - L_{avg}\right)^2.
\end{equation}
\listeq{Balanceamento local com carga acumulada}

Desta forma, as decisões tomadas nos \textit{sub-QUBOs} anteriores influenciam a
energia do \textit{sub-QUBO} atual através de termos lineares.

\section{Resolução Iterativa}
\label{sec:resolucao-iterativa-subqubo}

Cada \textit{sub-QUBO} é resolvido sequencialmente com o mesmo \textit{solver}
\acs{QAOA} usado no caso direto. Seja $a_i^{(t)}$ o núcleo atribuído à entidade
$e_i \in \mathcal{G}_t$ na iteração $t$. Depois de resolver um
\textit{sub-QUBO} viável, as atribuições são fixadas e a carga acumulada é
atualizada:

\begin{equation}
    \phi^{(t+1)}_k =
    \phi^{(t)}_k
    +
    \sum_{\substack{e_i \in \mathcal{G}_t \\ a_i^{(t)} = k}} w_i.
\end{equation}
\listeq{Atualização da carga acumulada após resolver um sub-QUBO}

O processo repete-se até todos os grupos terem sido resolvidos. Se algum
\textit{sub-QUBO} não produzir uma atribuição viável, a pipeline termina antes
de atualizar as atribuições finais e o vetor de carga, evitando propagar uma
solução inválida para as iterações seguintes.

No final, as atribuições parciais são combinadas numa única solução global. Esta
solução é validada reconstruindo a bitstring correspondente e avaliando a energia
na formulação \acs{QUBO} global. A decomposição em \textit{sub-QUBOs} não é
equivalente a resolver o \acs{QUBO} completo de uma só vez; trata-se de uma
aproximação iterativa que reduz a dimensão de cada problema resolvido, propagando
as decisões anteriores através do vetor $\boldsymbol{\phi}$.

\subsection{Exemplo de Extração dos Sub-QUBOs}

Retomando os quatro \textit{bundles} do exemplo anterior, as suas cargas
agregadas são:

\begin{equation}
    W_0=0.21,\qquad
    W_1=0.41,\qquad
    W_2=0.61,\qquad
    W_3=0.81.
\end{equation}

Como existem quatro \textit{bundles} e dois núcleos, o QUBO global possui
$4\times 2=8$ variáveis. Este valor excede o limite $q_{\max}=4$.

Cada \textit{sub-QUBO} pode incluir, no máximo,

\begin{equation}
    \left\lfloor\frac{q_{\max}}{K}\right\rfloor
    =
    \left\lfloor\frac{4}{2}\right\rfloor
    =2
\end{equation}

\textit{bundles}. Aplicando a ordenação \texttt{WEIGHT\_DESCENDING}, são formados
os seguintes grupos:

\begin{equation}
    \mathcal{G}_0=\{B_3,B_2\},
    \qquad
    \mathcal{G}_1=\{B_1,B_0\}.
\end{equation}

Esta divisão não cria novos \textit{bundles}. Cada grupo determina apenas
as variáveis de $Q^G$ que serão incluídas no respetivo \textit{sub-QUBO}.

\subsubsection{Primeiro sub-QUBO}

O primeiro grupo contém os \textit{bundles} $B_3$ e $B_2$. As suas
variáveis de atribuição são

\begin{equation}
    \mathcal{V}_0=
    \{x_{3,0},x_{3,1},x_{2,0},x_{2,1}\}.
\end{equation}

O primeiro \textit{sub-QUBO} corresponde à submatriz de $Q^G$ associada
a estas variáveis:

\begin{equation}
    Q^{(0)}=Q^G[\mathcal{V}_0,\mathcal{V}_0].
\end{equation}

Como o exemplo não inclui processos de tempo real, a carga acumulada
inicial é

\begin{equation}
    \boldsymbol{\phi}^{(0)}=(0,0).
\end{equation}

Suponha-se que a solução obtida atribui $B_3$
ao núcleo 0 e $B_2$ ao núcleo 1. A carga acumulada passa então a ser

\begin{equation}
    \boldsymbol{\phi}^{(1)}
    =
    (W_3,W_2)
    =
    (0.81,0.61).
\end{equation}

\subsubsection{Segundo sub-QUBO}

O segundo grupo contém os \textit{bundles} $B_1$ e $B_0$, sendo consideradas
as variáveis

\begin{equation}
    \mathcal{V}_1=
    \{x_{1,0},x_{1,1},x_{0,0},x_{0,1}\}.
\end{equation}

A submatriz correspondente é extraída de $Q^G$. Antes da sua resolução,
a carga acumulada no primeiro \textit{sub-QUBO} é incorporada nos
elementos diagonais. Os acréscimos aplicados são:

\begin{table}[H]
\centering
\begin{tabular}{|c|c|}
\hline
\textbf{Variável} & \textbf{Acréscimo diagonal} \\
\hline
$x_{1,0}$ & $2(0.41)(0.81)=0.6642$ \\
$x_{1,1}$ & $2(0.41)(0.61)=0.5002$ \\
$x_{0,0}$ & $2(0.21)(0.81)=0.3402$ \\
$x_{0,1}$ & $2(0.21)(0.61)=0.2562$ \\
\hline
\end{tabular}
\caption{Bias introduzido no segundo sub-QUBO pela carga acumulada.}
\label{tab:exemplo-bias-segundo-subqubo}
\end{table}

Como o núcleo 0 possui inicialmente uma carga superior, as variáveis que
representam novas atribuições a esse núcleo recebem um acréscimo maior.

Neste caso, a solução atribui $B_1$ ao núcleo 1 e $B_0$ ao núcleo 0. A
carga final é, portanto,

\begin{equation}
\begin{aligned}
    \phi^{(2)}_0 &= W_3+W_0=0.81+0.21=1.02,\\
    \phi^{(2)}_1 &= W_2+W_1=0.61+0.41=1.02.
\end{aligned}
\end{equation}

A solução final obtida é:

\begin{table}[H]
\centering
\begin{tabular}{|c|c|c|c|}
\hline
\textbf{Núcleo} & \textbf{Bundles} & \textbf{Processos} &
\textbf{Carga total} \\
\hline
0 & $B_3,B_0$ & $p_7,p_8,p_1,p_2$ & $1.02$ \\
1 & $B_2,B_1$ & $p_5,p_6,p_3,p_4$ & $1.02$ \\
\hline
\end{tabular}
\caption{Resultado final do exemplo de resolução por sub-QUBOs.}
\label{tab:resultado-exemplo-subqubos}
\end{table}